\section{Mapping}
\label{sec:mapping}

The mapping layer provides bidirectional conversion between the canonical
semantic model and each storage engine's internal representation.

\subsection{Canonical Mapping}
\label{sec:mapping:canonical}

The canonical mapping $\phi: \mathcal{F} \to \mathfrak{F}$ converts a
generic \texttt{SemanticFact} into a \texttt{CanonicalFact}, which is the
intermediate representation used for cross-engine comparison. The inverse
mapping $\psi: \mathfrak{F} \to \mathcal{F}$ reconstructs the original fact.

\[
\psi(\phi(f)) = f \quad \text{(lossless)}
\]

\subsection{KG Mapping}
\label{sec:mapping:kg}

The KG mapping $\phi_{\mathrm{KG}}$ decomposes a \texttt{SemanticFact}
into reified triples:

\begin{enumerate}
    \item Create an event node $e$ with type \texttt{rdf:Statement}.
    \item Add triple $(e, \texttt{rdf:subject}, f.\text{subject})$.
    \item Add triple $(e, \texttt{rdf:predicate}, f.\text{relation})$.
    \item For each object $\omega_j$, add triple
          $(e, \texttt{rdf:object}_j, \omega_j)$.
\end{enumerate}

Reconstruction requires $k$ joins to recover the original fact from
its constituent triples.

\subsection{Sheaf Mapping}
\label{sec:mapping:sheaf}

The Sheaf mapping $\phi_{\mathrm{Sheaf}}$ stores a \texttt{SemanticFact}
directly as a section in its assigned context. No decomposition occurs.
The fact is registered with the open set corresponding to its context and
with ancestor contexts via restriction maps.

\subsection{Equivalence Verification}
\label{sec:mapping:verification}

The \texttt{QueryEngine.results\_equivalent()} method compares result sets
from different engines by canonical fact ID equality. The benchmark
framework verifies equivalence across all queries before recording
performance metrics.
